% =====================================================================
%  Chapter 14 — Features from First Principles
% =====================================================================

\chapter{Features from First Principles}
\label{ch:features}

The model's input is twelve numbers per fund per month. This chapter defines
each of them, states why it is there, and derives the two that are not
immediate --- the Amihud illiquidity measure and the beta --- because a
feature that is merely named is a feature that has not been checked.

\section{The momentum family}

Three features measure the same underlying quantity at different lengths,
and the reason for having three rather than one is that the length is the
parameter most likely to be wrong.

\[
\begin{aligned}
\mathrm{mom}_{12,1,t} &= \frac{P_{t-1}}{P_{t-13}} - 1, \\
\mathrm{mom}_{6,1,t}  &= \frac{P_{t-1}}{P_{t-7}} - 1, \\
\mathrm{mom}_{36,t}    &= \frac{P_t}{P_{t-36}} - 1.
\end{aligned}
\]

The subscript ``1'' is the skip. The most recent month's return is
excluded, and the reason is worth stating carefully because it is the
single most consequential convention in the feature layer.

The argument is that a price that has just risen has already been
accompanied by the buying that pushed it there. The cross-sectional
empirical regularity is that the one-month reversal is large and negative
--- the short-horizon effect that Chapter 11 noted institutions cannot
scale --- and that including the most recent month in a momentum measure
contaminates it with exactly that effect. A twelve-month momentum that
includes the last month is a mixture of a genuine medium-horizon signal
and a strong short-horizon reversal, and the mixture's weight depends on
how much of the move happened in the last month.

The thirty-six-month feature has no skip, for a different reason. At that
length the one-month effect is a rounding error and the extra
complexity of a skip is not worth it.

\section{The reversal feature}

\[
\mathrm{rev}_{1,t} = -r_{t-1}.
\]

One number, and it is the negative of last month's return. This is the
short-horizon reversal effect, and it is included in the feature set for
two reasons. The first is that it is the strongest single regularity in
the cross-section of liquid ETFs, and a model that does not contain it
would be unable to express the fact that a fund that rose last month is
currently expensive in the sense that matters. The second is that the
book is testing a hypothesis, and the hypothesis includes the claim that
the reversal is *stronger* than the momentum at short horizons. Letting
the model see both, and weighting them by the posterior, is how that claim
gets tested rather than asserted.

\section{The volatility family}

\[
\mathrm{vol}_{12,t} = \sqrt{12}\cdot s_{12,t},
\]
where \(s_{12,t}\) is the sample standard deviation of the trailing
twelve monthly returns, and
\[
\mathrm{volRatio}_t = \frac{\mathrm{vol}_{6,t}}{\mathrm{vol}_{12,t}} - 1.
\]

The first is the level of risk, and it is a feature because a forward
return distribution conditional on recent volatility differs from the
unconditional one --- low-volatility assets have historically earned
slightly less than high-volatility assets, which is the low-beta and
minimum-volatility premium, and the model should be able to express that
view if the data supports it.

The second is the \emph{change} in risk, and it is included for a
different reason: it is a crude measure of the information state. A fund
whose recent volatility has risen relative to its own year is in a
different regime from one whose volatility has fallen, and a signal that
ignores that distinction is treating them as the same asset.

\section{The trend, drawdown and exposure features}

\[
\mathrm{trend}_{12,t} = \frac{1}{12}\sum_{j=0}^{11}
  \mathbb{1}\{r_{t-j} > 0\} - \frac12,
\]
the excess of positive months over negative ones in the trailing year,
centred. This is a distributional shape measure rather than a magnitude
measure: a fund that rose by ten per cent in one month and fell by ten per
cent in eleven has the same twelve-month return as one that rose in ten
months and fell in two, and the two are in different places in the
distribution.

\[
\mathrm{dd}_{12,t} = \frac{P_t}{\max_{j=0,\dots,12} P_{t-j}} - 1,
\]
the distance from the twelve-month high. This is the feature most closely
related to the behavioural literature of Part II, and it is the one that
speaks most directly to the drawdown brake of Chapter 28: a fund already
far below its high is in the state in which a human holder is most likely
to sell, and the model can be asked what it expects from that state.

The two market-exposure features are the standard regression coefficients
of the fund's monthly return on the benchmark's:

\[
\beta_{36,t} = \frac{\operatorname{Cov}(r_i, r^m)_{36}}{\operatorname{Var}(r^m)_{36}},
\qquad
\rho_{36,t} = \operatorname{Corr}(r_i, r^m)_{36},
\]
both estimated over the trailing thirty-six months.

A high beta is not, by itself, a prediction. It is a statement about how
much of the market's risk the fund carries, and its usefulness here is
that it lets the model distinguish a fund that is a levered version of the
market from one that is genuinely different. A low correlation with the
market is a statement about diversification potential, and it is the
feature that makes the cross-sectional model able to express a preference
for a fund that adds something to a portfolio that already holds the
market.

\section{The liquidity feature, derived}

The Amihud (2002) measure is the ratio of the absolute price change to
the dollar volume, averaged over a window. The book uses the monthly
form

\[
\mathrm{illiq}_{12,t} = \frac{1}{12}\sum_{j=0}^{11}
  \frac{|r_{t-j}|}{V_{t-j}\cdot P_{t-j}},
\]
scaled by \(10^6\) for legibility.

The intuition is that a dollar traded in a thin market moves the price
more, so the price change per unit of money is a measure of how expensive
it is to trade. A high value means the fund moves a lot for little money:
it is illiquid, and it is expensive to hold in the sense that a position
in it cannot be exited at the quoted price.

The economic content is that illiquidity is priced: illiquid assets earn
a premium. Whether that premium is available to a single investor is a
separate question, and the book's answer --- developed in Chapter 33 ---
is that the spread cost the investor actually pays is a significant
fraction of it.

\section{Standardisation, and why it is per-date}

Every feature is standardised within each date:
\[
z_{i,t} = \frac{x_{i,t} - \bar{x}_t}{s_t},
\]
where \(\bar{x}_t\) and \(s_t\) are the mean and standard deviation of
that feature across the funds that exist on date \(t\).

Two properties matter. The first is that the standardisation is computed
per date, so no statistic computed over the whole sample ever enters the
model; there is no in-sample mean to leak. The second is that the
standardisation is computed on the \emph{available} cross-section, which
changes over time, so the scale of the signal is comparable across dates
even as the universe grows.

It also means that the model is estimating a \emph{relative} quantity: the
predicted cross-sectional rank of a fund's forward return against the funds
that existed alongside it. That is the right target for a tilt that can
only move money within a sleeve, and it is the reason the feature layer
needs no level information at all.

\section{What the feature set is for}

The honest summary of this chapter is that the feature set is small,
price-based, and deliberately incapable of expressing a fundamental view.
It can express momentum, reversal, volatility, trend shape, drawdown
state, market exposure, diversification and liquidity --- and nothing else.

That is a limitation and it is also the point. A system with twelve
price-based features and a handful of configurations has a small enough
hypothesis space that the selection machinery of Part IV has a chance of
working. A system with twelve hundred features from a data vendor's
fundamentals database does not, and Chapter 33 will show what happens when
the trial family is small enough to be evaluated honestly.

\begin{keytakeaways}
\item Twelve features in four families: momentum (with a one-month skip),
reversal, volatility and its change, trend shape and drawdown state,
market beta and correlation, and Amihud illiquidity.
\item The one-month skip on momentum is not a convention; it is the removal
of a strong short-horizon reversal from a medium-horizon signal.
\item Standardisation is per date and on the available cross-section, so no
whole-sample statistic ever enters the model and no scale is carried
across dates.
\item The feature set can express relative preference within the funds that
exist alongside a fund, which is exactly what a within-sleeve tilt needs
and exactly what a level-based forecast would not.
\end{keytakeaways}
